\begin{tabularx}{\textwidth}{@{}ll>{\raggedright\arraybackslash}X>{\raggedright\arraybackslash}X@{}}
  \toprule
  \textbf{Herramienta} & \textbf{Versión} & \textbf{Para qué se usó} & \textbf{Cómo se verificó lo obtenido} \\
  \midrule
  Claude & Opus 5.5, Sonnet 5.5 & Buscar artículos candidatos. & Abrí cada DOI y comparé autores, año, título, revista, volumen y páginas. \\
  \addlinespace
  Claude & Opus 5.5, Sonnet 5.5 & Ubicar datos concretos en los artículos. & Los comprobé en el PDF de cada artículo (sección y página). \\
  \addlinespace
  Claude & Opus 5.5, Sonnet 5.5 & Explicar conceptos (DOI, biblatex, pull request). & Lo comprobé con las instrucciones de la actividad al aplicarlo en el repositorio. \\
  \addlinespace
  Claude & Opus 5.5, Sonnet 5.5 & Orientación sobre LaTeX, BibTeX y GitHub. & Obtuve el BibTeX desde doi.org y probé los pasos en GitHub. \\
  \addlinespace
  Claude & Opus 5.5, Sonnet 5.5 & Revisar ortografía y redacción de mis textos. & Leí cada corrección y la acepté o la cambié yo; las ideas son mías. \\
  \addlinespace
  Claude & Opus 5.5, Sonnet 5.5 & Contar palabras de los resúmenes y detectar errores de formato en mi archivo .tex. & Comprobé los resultados en mi propio archivo. \\
  \addlinespace
  Claude & Opus 5.5, Sonnet 5.5 & Sugerir ideas y estructura para las notas personales. & Reescribí las notas con mis palabras y comprobé que lo que afirmo del sistema fuera cierto. \\
  \bottomrule
\end{tabularx}
